\section{Discussion}

The benchmark does not support a simple conclusion that one rendering strategy is always better than the other. Instead, it points to a trade-off between user-facing load behaviour and total client-side energy use. Within this Next.js prototype, SSR reduced median total energy in all three scenarios and did so most clearly in the \textit{static} and \textit{massive} cases. CSR, on the other hand, consistently produced a smaller measured navigation response and in the \textit{dynamic} scenario produced lower LCP and FCP values. The overall pattern is therefore mixed: different metrics favour different strategies, even though the energy results lean toward SSR.

This pattern fits the design of the prototype. In the SSR variant, the main data transformation was completed before the page was delivered, whereas in the CSR variant the browser performed that work after hydration and after a follow-up inventory request. Because both variants were implemented within the same Next.js application and shared the same interface and data model, the comparison remained focused on where the computation took place rather than on framework-level differences.

The timing results help explain why CSR can still look attractive in some scenarios despite its higher energy cost. A smaller initial response and deferred computation can support better early paint metrics, especially when the first visual state is a relatively light shell and the expensive transformation is postponed until after hydration. That interpretation fits the \textit{dynamic} scenario, where CSR achieved lower LCP and FCP than SSR and also a smaller measured navigation response. Better early paint metrics do not necessarily mean less total client work, however. If the client still has to fetch the inventory and execute the heavy transformation locally, some of the computational cost is shifted to a later point in the page lifecycle rather than removed.

The three scenarios help qualify that trade-off. In the \textit{static} scenario, SSR still showed lower total energy even though the dataset was small, which suggests that the extra client round trip and hydration cost were not offset by a clear latency gain in this prototype. In the \textit{dynamic} scenario, the comparison remained the most mixed: energy still favoured SSR by median, but CSR showed the clearest paint advantage. In the \textit{massive} scenario, the cost of leaving the transformation in the browser became much more visible. SSR not only reduced median total energy there, but also showed a much tighter energy distribution and a lower median JavaScript heap sample than CSR. This is the clearest case in the dataset where browser-side workload appears to matter directly.

Several limitations should keep the interpretation cautious. First, energy was estimated from battery current and voltage polling via \texttt{adb dumpsys battery} at 0.25-second intervals rather than from dedicated hardware power-rail counters. This is a practical and defensible method for the available device, but it is still less precise than laboratory-grade power instrumentation. Second, the final conclusions rely on one battery-powered wireless run on a single Samsung Galaxy A53. The randomised order of conditions, fixed screen brightness, and narrow temperature range reduce confounding influences, but they do not remove device-specific variability. Third, the energy distributions contained clear high outliers, particularly in CSR \textit{massive} and CSR \textit{static}, which is why the study emphasises medians, IQRs, and non-parametric testing. Finally, earlier USB-connected runs were still useful for validating the tooling and the collection pipeline, but USB power could affect battery-based readings. For that reason, the direct energy interpretation in this paper is based on the corrected battery-powered wireless benchmark rather than on the earlier wired measurements.
